\section{Efficient Global Optimization:\ Proofs}\label{app:optimization}
%!TEX root = ../main.tex


\gatedComplexity*
\begin{proof}
	Applying the solution mapping from the proof of \cref{thm:unconstrained-equivalence} (see \cref{app:convex-forms}) we find that taking
	\[
		v_i' = W^{*(1)}_{\cdot, i} * W^{*(2)}_{i},
	\]
	for each \( D_i \in \calD \) yields a global minimizer of C-ReLU.
	Now we apply the iteration complexity of FISTA~\citep[Theorem 4.4]{beck2009fista} to obtain an \( \epsilon \)-optimal solution in
	\begin{align*}
		T & \leq \frac{\rbr{2 \lambda_{\text{max}}\rbr{M^\top M}\sum_{D_i \in \tilde \calD} \norm{v_i'}_2^2}^{1/2}}{\epsilon^{1/2}}            \\
		  & = \frac{\rbr{2 \lambda_{\text{max}}\rbr{M^\top M}\sum_{D_i \in \tilde \calD} \norm{W^{*(1)}_{\cdot, i}* W^{*(2)}_{i}*}_2^2}^{1/2}}{\epsilon^{1/2}},
	\end{align*}
	iterations.
	Note that we have used the fact that \( \lambda_{\text{max}}(M^\top M) \) is the Lipschitz smoothness constant of the squared-error loss.

\end{proof}

\approxDecompBounds*
\begin{proof}
	First-order optimality conditions for \( \tilde w \) imply
	\begin{align*}
		-\tilde X^\top \rbr{b - \tilde X \tilde w}_+
		 & \in \rho \cdot \partial \norm{\tilde w}_2.
	\end{align*}
	Noting that every vector in \( \partial \norm{\tilde w}_2 \) has norm at most
	\( 1 \), we deduce
	\begin{align*}
		\norm{\tilde X^\top \rbr{b - \tilde X \tilde w}_+}_2
		 & \leq \rho                                        \\
		\implies
		\norm{\rbr{b - \tilde X \tilde w}_+}_2
		 & \leq \frac{\rho}{\sigma_{\text{min}}(\tilde X)}  \\
		\iff \norm{\rbr{\max\cbr{- \tilde X u, 0} - \tilde X \tilde w}_+}_2
		 & \leq \frac{\rho}{\sigma_{\text{min}}(\tilde X)}.
		\intertext{since \( \tilde X \) is full row-rank and by definition of \( b \).
			Using the fact that only positive elements contribute to the norm, we obtain the following two inequalities:}
		\norm{\rbr{\tilde X \tilde w}_-}_2
		 & \leq \frac{\rho}{\sigma_{\text{min}}(\tilde X)}  \\
		\norm{\rbr{\tilde X (u + \tilde w)}_-}_2
		 & \leq \frac{\rho}{\sigma_{\text{min}}(\tilde X)},
	\end{align*}
	Recalling \( \tilde v = u + \tilde w \) and summing gives the first result.

	If \( \rho > 0 \), then it is easy to observe
	(i.e. by arguing via contradiction) that \( \tilde w \) must
	have smaller norm than any \( w' \) in a feasible decomposition \( (v', w') \).
	Thus, it must also have smaller norm than \( \bar w \) from the SOCP cone
	decomposition:
	\[ \norm{\tilde w}_2 \leq \norm{\bar w}_2. \]
	Since the proof of \cref{prop:socp-decomp} relies on controlling
	only the norm of \( \bar w \), the conclusion of that theorem also applies
	to \( (\tilde v, \tilde w) \).

	Finally, suppose \( X \) is not full row-rank. For \( \rho > 0 \),
	\cref{eq:cd-approx} is equivalent to solving
	\begin{equation}\label{eq:cd-approx-constr}
		\textbf{CD-A}: \, \min_{w} g(w, \rho) : =
		\half
		\|
		(b - \tilde X w)_+
		\|_2^2
		+ \rho P(w) \quad \text{s.t.} \norm{w}_2 \leq \norm{\bar w}_2,
	\end{equation}
	where \( \bar w \) is the norm of the minimum-norm (with respect to \( w \)
	only) solution to the cone decomposition problem.
	This is a minimization problem with compact constraint set;
	since \( g \) is continuous in both \( w \) and \( \rho \),
	we may apply Berge's maximum theorem~\citep{ausubel1993generalized}
	to obtain find
	\[
		g^*(\rho) = \min_{w}
		\cbr{ g(w, \rho) : \norm{w}_2 \leq \norm{\bar w}_2},
	\]
	is continuous.
	Since the cone decomposition is realizable at \( \rho = 0 \),
	\( g^*(0) = 0 \) and any sequence \( \rho_k \) converging to \( 0 \)
	satisfies
	\( \lim_k g^*(\lambda_k) = 0 \).

	Let \( \tilde w_k \) be the sequence of minimizers associated with \( \rho_k \).
	Since \( \tilde w_k \) is bounded, it has at least one convergent subsequence.
	Let \( \tilde w_0 \) be the associated limit point.
	Since \( g \) is continuous in \( w \) and \( \rho \), we find
	\[ g(w_0, 0) = \lim_k g(\tilde w_k, \rho_k) = \lim_k g^*(\rho_k) = 0, \]
	which shows that \( \tilde w_0 \) is a feasible decomposition.
	This completes the proof.
\end{proof}

\begin{restatable}{proposition}{approxDecompSubmatrices}\label{prop:approx-decomp-submatrices}
	Suppose \( \tilde w \) is a minimizer of \eqref{eq:cd-approx} and let \( \tilde v = u + \tilde w \).
	There exists \( \calJ \subseteq [n] \) such that
	\[
		\|(\tilde X \tilde w)_-\|_2 + \|(\tilde X \tilde v)_-\|_2
		\leq \frac{2\rho}{\sigma_{\text{min}}(\tilde X_{\calJ})},
	\]
	where \( \sigma_{\text{min}}(\tilde X_{calJ}) \) is the minimum
	(possibly zero) singular value of the sub-matrix \( \tilde X_{\calJ} \).
\end{restatable}
\begin{proof}
	First-order optimality conditions for \( \tilde w \) imply
	\begin{align*}
		-\tilde X^\top \rbr{b - \tilde X \tilde w}_+
		 & \in \rho \cdot \partial \norm{\tilde w}_2.
	\end{align*}
	Let \( \calJ = \cbr{i \in [n] : b_i - \abr{\tilde x_i, \tilde w} > 0 } \)
	and define \( B \) be a diagonal matrix such that \( B_{jj} = 1 \) if
	\( j \in \calJ \) and \( 0 \) otherwise.
	In this notation, the optimality conditions can be written as
	\begin{align*}
		-\tilde X^\top B\rbr{b - \tilde X \tilde w}
		 & \in \rho \cdot \partial \norm{\tilde w}_2  \\
		\iff
		-(B \tilde X^\top)\rbr{b - \tilde X \tilde w}
		 & \in \rho \cdot \partial \norm{\tilde w}_2.
	\end{align*}
	Noting that every vector in \( \partial \norm{\tilde w}_2 \) has norm at most
	\( 1 \), we deduce
	\begin{align*}
		\norm{(B \tilde X)^\top \rbr{b - \tilde X \tilde w}}_2
		 & \leq \rho                                                \\
		\iff
		\norm{\tilde X_{\calJ}^\top \rbr{b - \tilde X \tilde w}_{\calJ}}_2
		 & \leq \rho                                                \\
		\implies
		\norm{\rbr{b - \tilde X \tilde w}_{\calJ}}_2
		 & \leq \frac{\rho}{\sigma_{\text{min}}(\tilde X_{\calJ})}.
	\end{align*}
	Recalling the definition of \( \calJ \) and proceeding as in the proof of
	\cref{prop:approx-decomp-bounds} gives the desired result.
\end{proof}

\begin{remark}
	The bound given in \cref{prop:approx-decomp-submatrices} may be
	vacuous when \( \tilde X_{\calJ} \) is not full row-rank since it
	concerns the minimum singular value, rather than the minimum \emph{non-zero}
	singular value.
	In this respect, it is unlike the other results given so which have relied
	on the minimum non-zero singular value.
	However, it is worth reporting since this bound may be non-vacuous even when
	\( \tilde X \) is not full row-rank and only the asymptotic result in
	\cref{prop:approx-decomp-bounds} applies.
\end{remark}

\begin{restatable}{theorem}{alConvergenceRate}\label{thm:al-convergence-rate}
	Let \( \rbr{\gamma^*, \zeta^*} \) be the minimum-norm maximizer of the Lagrange dual of Problem~\ref{convex-forms:eq:convex-relu-mlp}.
	Assume \( \delta > 0 \) is fixed and at each iteration \cref{eq:al-subroutine} is carried out exactly.
	Then, the AL method computes an \( \epsilon \)-optimal estimate \( \rbr{\gamma_\epsilon, \zeta_\epsilon} \) in
	\begin{equation*}
		T \leq \frac{\norm{\gamma^*}_2^2 + \norm{\zeta^*}_2^2}{\delta \epsilon},
	\end{equation*}
\end{restatable}
\begin{proof}
	Let \( d \) be the Lagrange dual function associated the Problem~\ref{convex-forms:eq:convex-relu-mlp}.
	We will show that the desired iteration complexity follows from standard results in the optimization literature.

	Firstly, it is well-known that if the primal objective is a proper, closed, convex function, then one iteration of the AL method with penalty strength \( \delta > 0 \) is equivalent to the following proximal-point step on the dual problem:
	\begin{align*}
		\rbr{\gamma_{k+1}, \zeta_{k+1}} & = \argmax_{\gamma \geq 0, \zeta \geq 0} \cbr{d(\gamma, \zeta) - \frac{1}{2 \delta} \sbr{\norm{\gamma - \gamma_k}^2 \norm{\zeta - \zeta_k}^2} }.
	\end{align*}
	See \citet[Section 5.4.6]{bertsekas1997nonlinear} for a proof of this fact.

	Invoking \citet[Theorem 2.1]{guler1991convergence} implies that the AL method attains the following convergence rate for the dual parameters:
	\begin{equation}
		d(\gamma^*, \zeta^*) - d(\gamma_k, \zeta_k) \leq \frac{\norm{\gamma^* - \gamma_0}_2^2 + \norm{\zeta^* - \gamma_0}_2^2}{\delta k}.
	\end{equation}
	Choosing \( \gamma_0 = \zeta_0 = 0 \) and re-arranging this equation gives the desired iteration complexity.
\end{proof}

\subsection{Data Normalization}\label{app:data-normalization}

Recall that the proximal gradient update has the form,
\begin{align*}
	\xkk & = \argmin_{x} \cbr{f(\xk) + \abr{\grad(\xk), x - \xk} + \frac{1}{2 \etak}\norm{x - \xk}_2^2 + g(x)}                         \\
	     & = \argmin_{x} \cbr{\frac{1}{2\etak} \norm{x - \rbr{\xk - \etak \grad(\xk)}}_2^2 + g(x)}.
	\intertext{Taking \( g(x) \) to be the group \( \ell_1 \) penalty, we have}
	\xkk & = \argmin_{x} \cbr{\frac{1}{2\etak} \norm{x - \rbr{\xk - \etak \grad(\xk)}}_2^2 + \lambda \sum_{g \in \calG} \norm{x_g}_2},
\end{align*}
where \( \calG \) is the set of group indices.
Letting \( \x^{+} = \xk - \etak \grad(\xk) \), the update takes the form (see, e.g. \citet[Section 2.3]{sra2012optimization}),
\begin{align*}
	\sbr{\xkk}_g & = \rbr{1 - \frac{\lambda}{\norm{\x^{+}_g}_2}}_+ \x^{+}_g,
\end{align*}
which establishes our claim that the proximal step~\eqref{unconstrained:eq:line-search-cond} is a thresholding operator.

Thresholding operators are sensitive to rounding and other forms of numerical error.
Indeed, it is not hard to see that using a perturbed gradient \( \hat \grad(\xk) = \grad(\xk) + \epsilon \) can lead to groups dropping out of the model (or staying in the model) when they should remain non-zero.
Thus, it is important to reduce numerical error as much as possible by improving the condition of other operations, like computing \( \grad(\xk) \).
We can use data normalization to partially achieve this goal.

In the remainder of this section, we restrict ourselves to the C-GReLU problem with squared loss,
\[
	\min_{v} \; \half \norm{\sum_{D_i \in \tilde \calD} D_i X v_i - y}_2^2 + \lambda \sum_{D_i \in \tilde \calD} \norm{v_i}_2.
\]
The Hessian of the smooth component of this problem is \( \nabla^2 f(v) = M^\top M \), where \( M \) is the ``expanded'' data matrix \( M = \sbr{D_1 X \; D_2 X \; \ldots D_{|\tilde \calD|} X} \).
Let \( h \in \R^{d}, h_i = \norm{X_{\cdot,i}}_2 \) and \( H = \text{diag}(h) \).
That is, \( H \) is a diagonal matrix with the column-norms of \( X \) along the diagonal.
Finally, define the column-normalized version of \( X \) to be \( N = H^{-1} X \).

It is not hard to see that the diagonal elements of \( N^\top N \) are \( 1 \) by construction.
Applying a trace bound, we have
\begin{equation}
	\lambda_{\text{max}}(N^\top N) \leq \text{trace}\rbr{N^\top N} = d.
\end{equation}
Now, consider the normalized version of the expanded data matrix, \( \tilde N = \sbr{D_1 N \; D_2 N \ldots} \).
Recalling each \( D_i \) is a diagonal matrix whose elements are either \( 0 \) or \( 1 \), we have

\[ N^\top D_i^\top D_i N = N^\top D_i N^\top \preceq N^\top N, \]
for each \( D_i \in \tilde \calD \) and the diagonal elements of this matrix are bounded by \( 1 \).
We conclude that
\[
	\lambda_{\text{max}}(\tilde N^\top \tilde N) \leq \text{trace}\rbr{\tilde N^\top \tilde N} \leq d \cdot |\tilde \calD|.
\]
This establishes the claim in Section~\ref{sec:efficient-fista} that data normalization can be used to upper-bound the maximum eigenvalues of the Hessian.

Moving on to computation of the gradient, note that the Hessian \( \tilde N^\top \tilde N \) will be low-rank as long as \( |\tilde \calD| * d > n \).
In fact, this is nearly always the case since we typically choose \( \tilde \calD \) to be as large as possible.
Thus, although the condition number of \( \nabla^2 f(v) \) is not well-defined, it is possible to reduce the maximum expansion entailed by the Hessian via column normalization.
Observing \( \nabla f(v) = \nabla^2 f(v) v - \tilde N^\top y \), we may expect conditioning of the gradient computation to improve.
Finally, since the C-GReLU is a linear model, transforming the weights as \( v_i' = H^{-1} v_i \) after optimization can be used to project the model back into the original data space, ensuring that data normalization has no effects outside of optimization.
